\documentclass[../main.tex]{subfiles}
\begin{document}
\chapter{2021-2022学年微积分（一）（上）期末考试}

\section{单项选择题(每小题 3 分，共 18 分)}
\begin{enumerate}
    \item 设函数$f(x)$在$(-\infty,+\infty)$内单调有界，$\{x_{n}\}$是一数列，则下列命题正确的是（\qquad）.

          \begin{tasks}[label=\textbf{\Alph*.}, label-width=1.5em, column-sep=1em](2)
              \task 若$\{x_{n}\}$收敛，则$\{f(x_{n})\}$收敛
              \task 若$\{x_{n}\}$单调，则$\{f(x_{n})\}$收敛
              \task 若$\{f(x_{n})\}$收敛，则$\{x_{n}\}$收敛
              \task 若$\{f(x_{n})\}$单调，则$\{x_{n}\}$收敛
          \end{tasks}
          \vspace{1em}

    \item 函数$f(x)=\lim_{n\to\infty}\frac{x^{n}+2}{x^{n}+1}$的间断点及类型是（\qquad）.

          \begin{tasks}[label=\textbf{\Alph*.}, label-width=1.5em, column-sep=1em](2)
              \task $x=-1$是第二类间断点
              \task $x=1$是第二类间断点
              \task $x=\pm1$均是第一类间断点
              \task $x=\pm 1$均是第二类间断点
          \end{tasks}
          \vspace{1em}

    \item $x\to0^{+}$时，与$\sqrt{x}$等价的无穷小量是（\qquad）.

          \begin{tasks}[label=\textbf{\Alph*.}, label-width=1.5em, column-sep=1em](2)
              \task $1-\mathrm{e}^{\sqrt{x}}$
              \task $\sqrt{1+\sqrt{x}}-1$
              \task $\ln\frac{1+x}{1-\sqrt{x}}$
              \task $1-\cos{\sqrt{x}}$
          \end{tasks}
          \vspace{1em}

    \item 设函数$f(x)$在$x=0$处连续，下列命题错误的是（\qquad）.

          \begin{tasks}[label=\textbf{\Alph*.}, label-width=1.5em, column-sep=1em](1)
              \task 若$\lim_{x\to0}\frac{f(x)}{x}$存在，则$f(0)=0$
              \task 若$\lim_{x\to0}\frac{f(x)+f(-x)}{x}$存在，则$f(0)=0$
              \task 若$\lim_{x\to0}\frac{f(x)}{x}$存在，则$f'(0)$存在
              \task 若$\lim_{x\to0}\frac{f(x)-f(-x)}{x}$存在，则$f'(0)$存在
          \end{tasks}
          \vspace{1em}

    \item 曲线$y=x\ln\left(\mathrm{e}+\frac{1}{x}\right)(x>0)$的渐近线条数为（\qquad）.

          \begin{tasks}[label=\textbf{\Alph*.}, label-width=1.5em, column-sep=1em](2)
              \task 0
              \task 1
              \task 2
              \task 3
          \end{tasks}
          \vspace{1em}

    \item 设$F(x)=\int^{x+2\pi}_{x}\mathrm{e}^{\sin t}\sin t\,\mathrm{d}t$，则$F(x)$（\qquad）.

          \begin{tasks}[label=\textbf{\Alph*.}, label-width=1.5em, column-sep=1em](2)
              \task 为正常数
              \task 为负常数
              \task 恒为0
              \task 不为常数
          \end{tasks}
          \vspace{1em}
\end{enumerate}

\section{填空题(每小题 4 分，共 16 分)}
\begin{enumerate}
    \setcounter{enumi}{6}
    \item 曲线$\begin{cases}
                  x = \arctan t, \\
                  y=\ln\sqrt{1+t^{2}}
              \end{cases}$对应于$t=1$处的法线方程为$=\underline{\hspace{3cm}}$.
          \vspace{0.5em}

    \item 曲线$y=x\sin x+2\cos x\left(-\frac{\pi}{2}<x<2\pi\right)$的拐点是$\underline
              {\hspace{3cm}}$.
          \vspace{0.5em}

    \item 曲线$y=\ln\cos x\left(0\leq x\leq\frac{\pi}{6}\right)$的弧长为$\underline
              {\hspace{3cm}}$.
          \vspace{0.5em}

    \item $y=2^{x}$的Maclaurin展开式中$x^{n}$项的系数为$\underline{\hspace{3cm}
              }$.
\end{enumerate}

\section{计算题(每小题 7 分，共 42 分)}
\begin{enumerate}
    \setcounter{enumi}{10}
    \item 已知$\lim_{x\to1}\frac{ax^{2}+x-3}{x-1}=b$，求常数$a,b$的值.
          \vspace{12em}

    \item 设$f(x)$为连续函数，且满足$f(x)=x^{2}-x\cdot f(2)+2\int_{0}^{1}f(x
              )\,\mathrm{d}x$，求$f(x)$.
          \vspace{11em}

    \item 求极限$l=\lim_{n\to\infty}\sum_{i=0}^{n-1}\frac{1}{n+i}$.
          \vspace{11em}

    \item 计算定积分$I=\int_{0}^{1}x^{3}\sqrt{1-x^{2}}\,\mathrm{d}x$.
          \vspace{11em}

    \item 设函数$f(x)=
              \begin{cases}
                  \lambda \mathrm{e}^{-\lambda x}, & x>0,    \\
                  0,                               & x\leq0,
              \end{cases}(\lambda>0)$，计算$\int_{-\infty}^{+\infty}xf(x)\,\mathrm{d}x$.
          \vspace{12em}

    \item 求微分方程$xy'+y-\mathrm{e}^{x}=0$，$y(2)=1$的特解.
          \vspace{12em}
\end{enumerate}

\section{综合题(每小题 7 分，共 14 分)}
\begin{enumerate}
    \setcounter{enumi}{16}
    \item 已知$f(x)$在$(-\infty,+\infty)$上连续，$f(x)=(x+1)^{2}+2\int_{0}^{x}
              f(t)\,\mathrm{d}t$，求$f^{(n)}(0)$的值（$n\geq2$）.
          \vspace{10em}

    \item 设抛物线$y=ax^{2}+bx+c$过原点，当$0\leq x\leq 1$时，$y\geq0$，又该抛物线与直线$x
              =1$及$x$轴围成平面图形的面积为$\frac{1}{3}$，求$a,b,c$的值，使得该图形绕$x$轴旋转一周而成的旋转体体积$V$最小.
          \vspace{10em}
\end{enumerate}

\section{证明题(每小题 5 分，共 10 分)}
\begin{enumerate}
    \setcounter{enumi}{18}
    \item 证明方程$\ln x = \frac{x}{\mathrm{e}}-2021$在区间$(0,+\infty)$内只有两个不同的实根.
          \vspace{10em}

    \item 设$f''(x)$在$[0,2]$上连续，且$|f''(x)|\leq M$，$f(1)=0$. 证明：
          \[
              \left|\int_{0}^{2}f(x)\,\mathrm{d}x\right|\leq\frac{M}{3}.
          \]
          \vspace{10em}
\end{enumerate}

\chapter{2021-2022学年微积分（一）（上）期末考试参考答案}
\section{单项选择题 (每小题 3 分，共 18 分)}
\begin{enumerate}
    \item \textbf{Solution}. B.

          考虑函数$f(x)=
              \begin{cases}
                  0, & x<0,   \\
                  1, & x\geq0
              \end{cases}$，它在$(-\infty,+\infty)$内单调不减且有界，取数列$x_{n}=\frac{(-1)^{n}}{n}$，则$x
                  _{n}\to0$，

          但当$n$为偶数时，$f(x_{n})=1$；当$n$为奇数时，$f(x_{n})=0$，所以$\{f(x_{n})
              \}$发散，A选项错误.

          若$\{x_{n}\}$单调，由于$f(x)$单调有界，所以$\{f(x_{n})\}$亦单调有界，必收敛，B选项正确.

          考虑函数$f(x)\equiv0$，它在$(-\infty,+\infty)$内单调不减且有界，取数列$x_{n}
              =n$，则$\{x_{n}\}$发散，

          但$\{f(x_{n})\}\equiv0$单调不减且收敛，故C，D选项错误.

    \item \textbf{Solution}. C.

          当$|x|>1$时，$x^{n}\to\infty$，所以$f(x)=\lim_{n\to\infty}\frac{x^{n}+2}{x^{n}+1}
              =1$；

          当$|x|<1$时，$x^{n}\to0$，所以$f(x)=\lim_{n\to\infty}\frac{x^{n}+2}{x^{n}+1}
              =2$；

          当$x=1$时，$f(x)=\lim_{n\to\infty}\frac{1+2}{1+1}=\frac{3}{2}$；

          当$x=-1$时，$f(x)=\lim_{n\to\infty}\frac{(-1)^{n}+2}{(-1)^{n}+1}$不存在.

          所以
          \[
              f(x)=
              \begin{cases}
                  1,           & |x|>1, \\
                  2,           & |x|<1, \\
                  \frac{3}{2}, & x=1,   \\
                  \text{不存在},  & x=-1
              \end{cases},
          \]
          因此$x=\pm1$均为$f(x)$的跳跃间断点，即第一类间断点.

    \item \textbf{Solution}. C.

          当$x\to0^{+}$时，
          \[
              \begin{gathered}
                  1-\mathrm{e}^{\sqrt{x}}=-\left(\mathrm{e}^{\sqrt{x}}-1\right)\sim -\sqrt{x},\\
                  \sqrt{1+\sqrt{x}}-1=\left(1+\sqrt{x}\right)^{\frac{1}{2}}-1\sim\frac{1}{2}\sqrt{x},\\
                  \ln\frac{1+x}{1-\sqrt{x}}=\ln\left(1+x\right)-\ln\left(1-\sqrt{x}\right)\sim
                  x+\sqrt{x}\sim\sqrt{x},\\
                  1-\cos{\sqrt{x}}=1-\left(1-\frac{1}{2}x+o(x)\right)\sim\frac{1}{2}x.
              \end{gathered}
          \]

    \item \textbf{Solution}. D.

          若$\lim_{x\to0}\frac{f(x)}{x}$存在，由于$f(x)$在$x=0$处连续，$f(0)=\lim_{x\to0}
              f(x)=\lim_{x\to0}\frac{f(x)}{x}\cdot \lim_{x\to0}x=0$，A选项正确.

          则$\lim_{x\to0}\frac{f(x)-f(0)}{x-0}=\lim_{x\to0}\frac{f(x)}{x}=f'(0)$存在，C选项正确.

          同理，若$\lim_{x\to0}\frac{f(x)+f(-x)}{x}$存在，则$\lim_{x\to0}\left(f(x)+f
              (-x)\right)=2f(0)=0$，B选项正确.

          对于D选项，考虑$f(x)=|x|$，则
          \[
              \lim_{x\to0^-}\frac{f(x)-f(-x)}{x}=\lim_{x\to0^-}\frac{(-x)-(-x)}{x}= 0,
          \]
          \[
              \lim_{x\to0^+}\frac{f(x)-f(-x)}{x}=\lim_{x\to0^+}\frac{x-x}{x}=0,
          \]
          故$\lim_{x\to0}\frac{f(x)-f(-x)}{x}=0$存在，但显然$f(x)$在$x=0$处不可导.

    \item \textbf{Solution}. B.

          当$x\to0^{+}$时，
          \[
              \lim_{x\to0^+}x\ln\left(\mathrm{e}+\frac{1}{x}\right)=\lim
              _{t\to+\infty}\frac{\ln(\mathrm{e}+t)}{t}=0,
          \]
          所以曲线$y=x\ln\left(\mathrm{e}+\frac{1}{x}\right)$没有竖直渐近线.

          当$x\to+\infty$时，
          \[
              \lim_{x\to+\infty}\frac{x\ln\left(\mathrm{e}+\frac{1}{x}\right)}{x}
              =1,\quad\lim_{x\to+\infty}\left[x\ln\left(\mathrm{e}+\frac{1}{x}\right)-x\right
              ]=\lim_{t\to0^+}\frac{\ln(\mathrm{e}+t)-1}{t}=\frac{1}{\mathrm{e}},
          \]
          所以曲线$y=x\ln\left(\mathrm{e}+\frac{1}{x}\right)$有斜渐近线$y=x+\frac{1}{\mathrm{e}}$.

          综上，曲线$y=x\ln\left(\mathrm{e}+\frac{1}{x}\right)$的渐近线条数为1.

    \item \textbf{Solution}. A.

          因被积函数$\mathrm{e}^{\sin t}\sin t$是周期为$2\pi$的函数，所以$F(x)=\int^{x+2\pi}
              _{x}\mathrm{e}^{\sin t}\sin t\,\mathrm{d}t=\int_{0}^{2\pi}\mathrm{e}^{\sin t}
              \sin t\,\mathrm{d}t$恒为常数.

          计算
          \[
              \int_{0}^{2\pi}\mathrm{e}^{\sin t}\sin t\,\mathrm{d}t =\int_{0}^{\pi}\mathrm{e}
              ^{\sin t}\sin t\,\mathrm{d}t+\int_{\pi}^{2\pi}\mathrm{e}^{\sin t}\sin t\,\mathrm{d}
              t.
          \]
          对第二项作代换$u=2\pi-t$，则
          \[
              \int_{\pi}^{2\pi}\mathrm{e}^{\sin t}\sin t\,\mathrm{d}
              t=\int_{\pi}^{0}\mathrm{e}^{\sin(2\pi-u)}\sin(2\pi-u)(-\mathrm{d}u)=-\int_{0}
              ^{\pi}\mathrm{e}^{-\sin u}\sin u\,\mathrm{d}u.
          \]
          所以
          \[
              \int_{0}^{2\pi}\mathrm{e}^{\sin t}\sin t\,\mathrm{d}t = \int_{0}^{\pi}\left
              (\mathrm{e}^{\sin t}-\mathrm{e}^{-\sin t}\right)\sin t\,\mathrm{d}t.
          \]
          因为在$(0,\pi)$上，$\sin t>0$，且$\mathrm{e}^{\sin t}>\mathrm{e}^{-\sin t}$，故
          \[
              \int_{0}^{\pi}\left(\mathrm{e}^{\sin t}-\mathrm{e}^{-\sin t}\right)\sin t\,\mathrm{d}
              t=\int_{0}^{2\pi}\mathrm{e}^{\sin t}\sin t\,\mathrm{d}t>0.
          \]
\end{enumerate}

\section{填空题(每小题 4 分，共 16 分)}
\begin{enumerate}
    \setcounter{enumi}{6}
    \item \textbf{Solution}. $y+x-\frac{1}{2}\ln 2-\frac{\pi}{4}=0$.

          当$t=1$时，$x=\arctan 1=\frac{\pi}{4}$，$y=\ln\sqrt{2}=\frac{1}{2}\ln 2$，且
          \[
              \frac{\mathrm{d}y}{\mathrm{d}x}=\frac{\frac{\mathrm{d}y}{\mathrm{d}t}}{\frac{\mathrm{d}x}{\mathrm{d}t}}
              =\frac{\frac{2t}{2(1+t^{2})}}{\frac{1}{1+t^{2}}}=t,
          \]
          所以$\left.\frac{\mathrm{d}y}{\mathrm{d}x}\right|_{t=1}=1$. 切线斜率为1，故法线斜率为$-1$，法线方程为$y-\frac{1}{2}\ln 2=-1\cdot\left(x
              -\frac{\pi}{4}\right)$，即
          \[
              y+x-\frac{1}{2}\ln 2-\frac{\pi}{4}=0.
          \]

    \item \textbf{Solution}. $(\pi,-2)$.

          计算
          \[
              y'=\sin x+x\cos x-2\sin x,\quad y''=\cos x+\cos x-x\sin x-2\cos x=-x\sin x,
          \]
          在$\left(-\frac{\pi}{2},2\pi\right)$内，令$y''=0$得$x=0$或$x=\pi$.

          当$x \in \left(-\frac{\pi}{2}, 0\right)\cup(0, \pi)$时，$y'' \leq 0$，$y''$在$x
              =0$两侧不变号，故$x=0$不是拐点；

          当$x \in (\pi, 2\pi)$ 时，$y'' > 0$，$y''$ 在 $x=\pi$ 两侧变号.

          故曲线的拐点为 $(\pi, -2)$.

    \item \textbf{Solution}. $\frac{1}{2}\ln 3$.

          因为$y'=\frac{-\sin x}{\cos x}=-\tan x$，所以$\mathrm{d}s=\sqrt{1+y'^{2}}\,\mathrm{d}
              x=\sqrt{1+\tan^{2}x}\,\mathrm{d}x=\sec x\,\mathrm{d}x$，

          故
          \[
              s=\int_{0}^{\frac{\pi}{6}}\sec x\,\mathrm{d}x=\ln\left|\sec x+\tan x\right
              |\bigg|_{0}^{\frac{\pi}{6}}=\ln\left(\frac{2}{\sqrt{3}}+\frac{1}{\sqrt{3}}\right
              )=\frac{1}{2}\ln 3.
          \]

    \item \textbf{Solution}. $\frac{(\ln 2)^{n}}{n!}$.

          因为$y=2^{x}=\mathrm{e}^{x\ln 2}=\sum_{n=0}^{\infty}\frac{(x\ln 2)^{n}}{n!}$，所以$x
                  ^{n}$项的系数为$\frac{(\ln 2)^{n}}{n!}$.
\end{enumerate}
\section{计算题(每小题 7 分，共 42 分)}
\begin{enumerate}
    \setcounter{enumi}{10}
    \item \textbf{Solution}. 由题可知$\lim_{x\to1}\left(ax^{2}+x-3\right)=a+
              1-3=0$，所以$a=2$.

          $b=\lim_{x\to1}\frac{2x^{2}+x-3}{x-1}=\lim_{x\to1}\frac{(2x+3)(x-1)}{x-1}=5$.

    \item \textbf{Solution}. 记$\int_{0}^{1}f(x)\,\mathrm{d}x=C$.

          令$x=2$，得$f(2)=4-2f(2)+2C$，所以$f(2)=\frac{4}{3}+ \frac{2}{3}C$.

          方程$f(x)=x^{2}-x\cdot f(2)+2\int_{0}^{1}f(x)\,\mathrm{d}x$两边从$0$积分到$1$得
          \[
              \begin{aligned}
                  C = {} & \int_{0}^{1}\left(x^{2}-x\cdot f(2)+2C\right)\,\mathrm{d}x                                       \\
                  = {}   & \frac{1}{3}-\frac{1}{2}f(2)+2C = \frac{1}{3}-\frac{1}{2}\left(\frac{4}{3}+\frac{2}{3}C\right)+2C \\
                  = {}   & \frac{5}{3}C-\frac{1}{3}.
              \end{aligned}
          \]
          解得$C=\int_{0}^{1}f(x)\,\mathrm{d}x=\frac{1}{2}$，所以
          \[
              f(x)=x^{2}-x\cdot f(2)+2\int_{0}^{1}f(x)\,\mathrm{d}x=x^{2}-\frac{5}{3}x+1.
          \]

    \item \textbf{Solution}. 由定积分的定义，
          \[
              \begin{aligned}
                  l = {} & \lim_{n\to\infty}\sum_{i=0}^{n-1}\frac{1}{n+i}= \lim_{n\to\infty}\frac{1}{n}\sum_{i=0}^{n-1}\frac{1}{1+\frac{i}{n}} \\
                  = {}   & \int_{0}^{1}\frac{1}{1+x}\,\mathrm{d}x = \ln 2.
              \end{aligned}
          \]

    \item \textbf{Solution}. 令$x=\sin\theta$，则$\mathrm{d}x=\cos\theta\,\mathrm{d}
              \theta$，
          \[
              \begin{aligned}
                  I = \int_{0}^{1}x^{3}\sqrt{1-x^2}\,\mathrm{d}x = {} & \int_{0}^{\frac{\pi}{2}}\sin^{3}\theta\sqrt{1-\sin^2\theta}\cos\theta\,\mathrm{d}\theta                                                             \\
                  = {}                                                & \int_{0}^{\frac{\pi}{2}}\sin^{3}\theta\cos^{2}\theta\,\mathrm{d}\theta = \int_{0}^{\frac{\pi}{2}}\sin^{3}\theta(1-\sin^{2}\theta)\,\mathrm{d}\theta \\
                  = {}                                                & \int_{0}^{\frac{\pi}{2}}\sin^{3}\theta\,\mathrm{d}\theta - \int_{0}^{\frac{\pi}{2}}\sin^{5}\theta\,\mathrm{d}\theta                                 \\
                  = {}                                                & \frac{2}{3}-\frac{4}{5}\cdot\frac{2}{3}= \frac{2}{15}.
              \end{aligned}
          \]

    \item \textbf{Solution}.
          \[
              \begin{aligned}
                  \int_{-\infty}^{+\infty}xf(x)\,\mathrm{d}x = {} & \int_{-\infty}^{0}xf(x)\,\mathrm{d}x+\int_{0}^{+\infty}xf(x)\,\mathrm{d}x                                                \\
                  = {}                                            & \int_{0}^{+\infty}\lambda x \mathrm{e}^{-\lambda x}\,\mathrm{d}x =-\int_{0}^{+\infty}x \mathrm{d}\mathrm{e}^{-\lambda x} \\
                  = {}                                            & -x\mathrm{e}^{-\lambda x}\bigg|_{0}^{+\infty}+\int_{0}^{+\infty}\mathrm{e}^{-\lambda x}\,\mathrm{d}x                     \\
                  = {}                                            & \int_{0}^{+\infty}\mathrm{e}^{-\lambda x}\,\mathrm{d}x = \frac{1}{\lambda}.
              \end{aligned}
          \]

    \item \textbf{Solution}. 方程变形为$y'+\frac{1}{x}y=\frac{\mathrm{e}^{x}}{x}$. 由一阶非齐次线性微分方程的通解公式得
          \[
              \begin{aligned}
                  y = {} & \mathrm{e}^{-\int\frac{1}{x}\,\mathrm{d}x}\left(C+\int \frac{\mathrm{e}^x}{x}\mathrm{e}^{\int\frac{1}{x}\,\mathrm{d}x}\,\mathrm{d}x\right) \\
                  = {}   & \frac{1}{x}\left(C+\int \mathrm{e}^{x}\,\mathrm{d}x\right) = \frac{1}{x}\left(C+\mathrm{e}^{x}\right).
              \end{aligned}
          \]
          将初始条件$y(2)=1$代入上式得$C=2-\mathrm{e}^{2}$，所以特解为$y=\frac{1}{x}\left
              (2-\mathrm{e}^{2}+\mathrm{e}^{x}\right)$.
\end{enumerate}

\section{综合题(每小题 7 分，共 14 分)}
\begin{enumerate}
    \setcounter{enumi}{16}
    \item \textbf{Solution}. 由于$f(x)$连续，所以其积分变上限函数$\int_{0}^{x}
              f(t)\,\mathrm{d}t$可导，

          从而方程$f(x)=(x+1)^{2}+2\int_{0}^{x}f(t)\,\mathrm{d}t$两边可导.两边求导得$f'
              (x)-2f(x)=2(x+1)$.

          由一阶非齐次线性微分方程的通解公式得
          \[
              \begin{aligned}
                  y = {} & \mathrm{e}^{-\int\left(-2\right)\,\mathrm{d}x}\left(C+\int 2(x+1)\mathrm{e}^{\int\left(-2\right)\,\mathrm{d}x}\,\mathrm{d}x\right)                                                \\
                  = {}   & \mathrm{e}^{2x}\left(C+2\int (x+1)\mathrm{e}^{-2x}\,\mathrm{d}x\right) = \mathrm{e}^{2x}\left[C+2\int x\mathrm{e}^{-2x}\,\mathrm{d}x+2\int \mathrm{e}^{-2x}\,\mathrm{d}x\right]   \\
                  = {}   & \mathrm{e}^{2x}\left[C+2\int x\mathrm{e}^{-2x}\,\mathrm{d}x-\mathrm{e}^{-2x}\right] = \mathrm{e}^{2x}\left[C-\int x\,\mathrm{d}\mathrm{e}^{-2x}-\mathrm{e}^{-2x}\right]           \\
                  = {}   & \mathrm{e}^{2x}\left[C-x\mathrm{e}^{-2x}+\int \mathrm{e}^{-2x}\,\mathrm{d}x-\mathrm{e}^{-2x}\right] = \mathrm{e}^{2x}\left[C-x\mathrm{e}^{-2x}-\frac{3}{2}\mathrm{e}^{-2x}\right] \\
                  = {}   & C\mathrm{e}^{2x}-x-\frac{3}{2}.
              \end{aligned}
          \]
          方程$f(x)=(x+1)^{2}+2\int_{0}^{x}f(t)\,\mathrm{d}t$两边令$x=0$得$f(0)=1$，

          所以$C=1+\frac{3}{2}=\frac{5}{2}$，$f(x)=\frac{5}{2}\mathrm{e}^{2x}-x-\frac{3}{2}$.当$n\geq2$时，
          \[
              f^{(n)}(0)=\left.\left(\frac{5}{2}\mathrm{e}^{2x}\right)^{(n)}
              \right|_{x=0}=5\cdot 2^{n-1}\mathrm{e}^{2x}\bigg|_{x=0}=5\cdot 2^{n-1}.
          \]

    \item \textbf{Solution}. 由抛物线过原点得$c=0$，且有
          \[
              \frac{1}{3}=\int_{0}^{1}\left(ax^{2}+bx\right)\,\mathrm{d}x=\frac{a}{3}+\frac{b}{2}
              ,
          \]
          即$b=\frac{2}{3}(1-a)$. 旋转体的体积
          \[
              \begin{aligned}
                  V = {} & \int_{0}^{1}\pi f^{2}(x)\,\mathrm{d}x = \pi\int_{0}^{1}\left(ax^{2}+bx\right)^{2}\,\mathrm{d}x = \pi\int_{0}^{1}x^{2}(ax+b)^{2}\,\mathrm{d}x   \\
                  = {}   & \pi\int_{0}^{1}(a^{2}x^{4}+2abx^{3}+b^{2}x^{2})\,\mathrm{d}x = \pi\left[\frac{a^2}{5}x^{5}+\frac{ab}{2}x^{4}+\frac{b^2}{3}x^{3}\right]_{0}^{1} \\
                  = {}   & \pi\left(\frac{a^2}{5}+\frac{ab}{2}+\frac{b^2}{3}\right) = \pi\left[\frac{a^2}{5}+\frac{a(1-a)}{3}+\frac{4(1-a)^2}{27}\right].
              \end{aligned}
          \]
          由$\frac{\mathrm{d}V}{\mathrm{d}a}=\pi\left[\frac{2}{5}a+\frac{1-2a}{3}-\frac{8(1-a)}{27}
                  \right]=0$解得$a=-\frac{5}{4}$，又
          \[
              \left.\frac{\mathrm{d}^{2}V}{\mathrm{d}a^{2}}
              \right|_{a=-\frac{5}{4}}=\pi\left(\frac{2}{5}-\frac{2}{3}+\frac{8}{27}\right
              )=\frac{4\pi}{135}>0,
          \]
          所以当$a=-\frac{5}{4}$，$b=\frac{3}{2}$，$c=0$时，旋转体体积$V$最小.
\end{enumerate}

\section{证明题(每小题 5 分，共 10 分)}
\begin{enumerate}
    \setcounter{enumi}{18}
    \item \textbf{Proof}. 令$f(x)=\ln x-\frac{x}{\mathrm{e}}+2021$，则$f'(x)
              =\frac{1}{x}-\frac{1}{\mathrm{e}}$，$f(x)$具有唯一驻点$x=\mathrm{e}$，

          且$f(x)$在$(0,\mathrm{e})$上单调递增，在$(\mathrm{e},+\infty)$上单调递减.

          因$f(0^{+})=-\infty<0$，$f(\mathrm{e})=2021>0$，$f(+\infty)=-\infty<0$，

          由介值定理和$f(x)$的单调性易知方程$\ln x = \frac{x}{\mathrm{e}}-2021$在区间$(
              0,+\infty)$内只有两个不同的实根.

    \item \textbf{Solution}. 由Taylor公式，将$f(x)$在$x=1$处展开得
          \[
              \begin{gathered}
                  f(x)=f(1)+f'(1)(x-1)+\frac{1}{2}f''(\xi)(x-1)^2=f'(1)(x-1)+\frac{1}{2}f''(\xi)(x-1)^2,
              \end{gathered}
          \]
          其中$\xi$介于$1$和$x$之间. 所以
          \[
              \begin{aligned}
                  \left|\int_{0}^{2}f(x)\,\mathrm{d}x\right| = {} & \left|\int_{0}^{2}\left[f'(1)(x-1)+\frac{1}{2}f''(\xi)(x-1)^{2}\right]\mathrm{d}x\right| \\
                  \leq  {}                                        & \frac{1}{2}\int_{0}^{2}|f''(\xi)|(x-1)^{2}\,\mathrm{d}x                                  \\
                  \leq  {}                                        & \frac{M}{2}\int_{0}^{2}(x-1)^{2}\,\mathrm{d}x = \frac{M}{3}.
              \end{aligned}
          \]
\end{enumerate}
\end{document}




